\documentclass[10pt]{article}
\usepackage[margin=0.8in]{geometry}
\usepackage{amsmath,amssymb,hyperref}
\hypersetup{colorlinks=true,urlcolor=blue,linkcolor=blue}
\setlength{\parindent}{0pt}\setlength{\parskip}{5pt}
\begin{document}
\begin{center}
{\Large\bfseries Problem 00000003901: a genuine Frobenius quotient length}\\[5pt]
lizaixi01 \quad AI-led research and formalization \quad 23 October 2026
\end{center}

\textbf{Result.} In the actual one-dimensional ring
\[
 R=\mathbb F_2[[X]],\qquad I=(X),\qquad q=2^e,
\]
the natural Frobenius quotient has actual $R$-module length
\[
 \ell_R\bigl(I/I^{[q]}\bigr)=q-1\qquad\text{for every }e\ge0.
\]
It is unbounded along $e$ and is not $O(q^{1-1})=O(1)$. In fact its normalization by $q^{1-1}$ cannot tend to any finite real coefficient. This refutes the necessary proposed exponent $n-1$ for the natural corrected quotient, without relying on finitely many sample values.

\textbf{Printed quotient direction.} The source writes $I^{[q]}/I$. A module quotient of ideals requires the denominator to be contained in the numerator. In our example $I^{[q]}\subseteq I$; Lean also proves $I\not\subseteq I^{[2]}$. Thus the printed direction is not even defined as an ordinary quotient already at $q=2$. Our substantive counterexample addresses the natural correction $I/I^{[q]}$, as did the earlier review's intended length discussion. The direction error is disclosed rather than used as the only objection.

\section*{1. Actual ring, ideal and Frobenius power}
The ring is the genuine formal power-series ring over \texttt{ZMod 2}, not an assigned object with a dimension label. Its discrete-valuation-ring structure, maximal ideal $I=(X)$ and actual Krull dimension $1$ are proved using Mathlib's power-series, DVR and PID dimension theorems. The source quantifies over every ideal of an $n$-dimensional ring and imposes no polynomial-ring-only restriction, so this ring is within that domain.

For an ideal $J$ in characteristic $p$, the standard $p^e$-th Frobenius power is generated by the $p^e$-th powers of its elements; see \cite{HK}, Introduction, printed p.~1. Lean uses exactly the ideal map under the actual iterated Frobenius endomorphism and proves
\[
 I^{[2^e]}=\operatorname{Ideal.map}(F^e)(I)=(X^{2^e}).
\]
It also proves this ideal is contained in $I$. The quotient denominator is its actual inverse image under the inclusion $I\hookrightarrow R$, so the resulting type is a genuine quotient $R$-module of $I$.

The example is a one-variable monomial ideal in a power-series ring. Its exponent set is $\{1,2,\ldots\}$ and its usual Newton polyhedron is $[1,\infty)$. No independently formalized Newton boundary-measure formula is needed or asserted: failure of the necessary exponent already precludes the full claimed exponent-and-coefficient statement.

\section*{2. A true linear equivalence computes the length}
Multiplication by $X$ is an actual $R$-linear bijection $R\to I$. Injectivity follows from $X\ne0$ and the domain property; surjectivity is the literal divisibility characterization of membership in $(X)$. For each $k\ge0$, its inverse image of $(X^{k+1})$ is $(X^k)$, by cancellation of $X$. Thus it induces the real $R$-linear equivalence
\[
 R/(X^k)\ \cong_R\ I/(X^{k+1}).
\]
Mathlib's DVR quotient-length theorem gives $\ell_R(R/\mathfrak m^k)=k$. Rewriting $\mathfrak m=(X)$ and transporting \emph{module length} through this linear equivalence yields
\[
 \ell_R(I/(X^{k+1}))=k.
\]
For $q=2^e>0$, take $k=q-1$. This proves the displayed all-$e$ formula for $I/I^{[q]}$. In particular, we never replace this module with $R/I^{[q]}$, whose length would be $q$, nor rename a vector-space dimension as an $R$-module length.

\newpage
\section*{3. The full asymptotic obstruction}
Lean's \texttt{Module.length} takes values in the extended natural numbers. The proof first establishes the actual finite value $q-1$ and its inequality to infinity. Only then does it extract the real numerical length $L(e)$ using \texttt{toNat}, proving $L(e)=(2^e-1)$.

Since $2^e-1\ge e$, for every real $C$ and natural cutoff $N$ one can choose
\[
 e=N+\lceil C\rceil_+ +2\ge N
 \quad\text{with}\quad C<L(e).
\]
This is proved with all constants and cutoffs quantified. It contradicts any eventual real bound, and hence the genuine Mathlib statement $L=O(1)$ is false. As $n=\dim R=1$, this is precisely the proposed $q^{n-1}$ scale. The source's choice of leading coefficient cannot repair the exponent: for every finite real $c$, the actual normalized function $L(e)/(2^e)^{1-1}$ does not tend to $c$; a finite limit would eventually bound it by $c+1$, contradicting the same cofinal inequality.

The formal result is this failure of the proposed zero exponent and finite leading coefficient. We do not additionally claim to have formalized a general theory of leading exponents, the exact leading exponent $1$, or Newton boundary measures. No finite numerical table substitutes for the asymptotic argument.

\section*{4. Earlier submission and verification}
Public \href{https://github.com/The-Last-Math-Competition/The-Last-Math-Competition/pull/285}{PR 285} already suggested a maximal-ideal Frobenius quotient as the obstruction, using a two-variable polynomial ring and the proposed length $q^2-1$. That exposed direction is attributed. The \href{https://github.com/The-Last-Math-Competition/The-Last-Math-Competition/pull/285#issuecomment-5965698008}{maintainer review} rejected a Lean project containing only four numeric length values and comparisons: no actual ring, ideal, Frobenius power, quotient module or length was represented, and no all-$q$ growth argument was proved.

This submission uses a simpler one-dimensional power-series DVR and formally closes every one of those bridges. The ring and Krull dimension, true ideal maps, real quotient submodules, actual linear equivalence, genuine \texttt{Module.length}, finite-value extraction and all-cutoff asymptotic contradiction are present in Lean. The change of example is explicit; it is not presented as a reproof of the old two-variable formula or as a proof about the reversed printed quotient.

All 32 named local declarations are audited, including the proof-bearing characteristic-two instance and linear equivalences. The package includes pinned Lean/Mathlib versions, a newly compiled root module, warning-as-error axiom audit, full fresh kernel replay and source-bound receipts. Full original bilingual source, LaTeX source and PDF are provided. There is no \texttt{sorry}, \texttt{native\_decide} or additional axiom.

\textbf{AI disclosure.} Research, Lean proofs and documents were produced by a Codex worker with separate AI supervision of source semantics, proof scope and PDF layout. Organizer acceptance and human verification are not asserted.

\begin{thebibliography}{1}
\bibitem{HK} C.~Huneke, M.~Katzman, R.~Y.~Sharp and Y.~Yao, \emph{Frobenius test exponents for parameter ideals in generalized Cohen--Macaulay local rings}, J. Algebra \textbf{305}(1) (2006), 516--539. Author version, Introduction p.~1 (PDF page~2, after repository cover), \href{https://eprints.whiterose.ac.uk/id/eprint/10176/1/Katzman5.pdf}{White Rose repository}; \href{https://doi.org/10.1016/j.jalgebra.2006.06.036}{doi:10.1016/j.jalgebra.2006.06.036}.
\end{thebibliography}
\end{document}
